\section{A signed horizon removes a quartic variable}
\label{sec:quartic}

The old direct source imposed a positive quadratic horizon such as
$t=16+4z^2$ or $t=128+4z^2$. This guaranteed positivity at the cost of
a coordinate that could not be removed affinely. The power machinery
already contains enough information to recover the sign. We use
that information and replace the square progression by a signed
affine progression.

\subsection{A uniform signed-horizon lemma}

Put $\ell=2^{20}$. Fix an integer $m\ge8$ divisible by four and impose
\[
 t=m(2z+1),\qquad z\in\Z.
\]
Then $t\ne0$, $|t|\ge m$ and $4\mid t$. Retain the source's signed
Pell frame, its reused even CRT index $C$, all orientation rows and
all actual power requests. Two growth choices will be used:
\begin{align}
 K&=\ell t, \label{eq:linear-growth}\\
 U&=a t^2,\qquad K=bU^2
 \quad(a,b\in\Z_{>0},\;ba^2m^3\ge\ell).
 \label{eq:quartic-growth}
\end{align}
The retained Pell request expresses $F=t^K$. This exponentiation
is part of the \emph{proof interpretation}, not the final polynomial.

\begin{lemma}\label{lem:signed-horizon}
At every integer zero of either modified source, $t\ge m$ and
$K\ge\ell t$. Every finite C222 refutation can be represented using
the new positive branch of the horizon progression.
\end{lemma}

\begin{proof}
The source's preliminary signed anchor argument uses squares and
absolute values. It recovers the canonical signed Pell pairs before
decoding arithmetic signs. Its actual frame includes the squares of
all requested bases, outputs and indices, in particular the changed
index $K$. Thus it still gives
\[
 W\ge258,\qquad |A|>W^W.
\]
For the affine index this lower bound on $W$ already follows from
$K^2\ge\ell^2m^2$; for the quartic index it follows from
$K=ba^2t^4$. The signed power lemma therefore applies at a nonzero
base. At a base of absolute value at least two, an integral exact
signed-index power must have nonnegative exponent.

For \eqref{eq:linear-growth}, $|t|\ge8$ makes the first power
force $K\ge0$. Since $K=\ell t$ and $t\ne0$, this proves $t\ge m$
immediately.

For \eqref{eq:quartic-growth}, the index $K=ba^2t^4$ is positive
and even before the sign of $t$ is known. The first power gives
$F=t^K>0$, with $F$ even and $F>|t|$. The retained CRT row is
\[
 Bj=C+F+1,\qquad C=2c.
\]
Its right side is odd, so $B$ is odd and nonzero. The self-radix
multiplier has the form $H=2+4S$ for an integer $S$, hence
$|BH|\ge2$. The exact power at that base gives
$x\ge0$ and $q=(BH)^x$.

It remains to rule out $x=0$. The beta endpoint rows, which were
never deleted, give
\[
 y=F(2v+1),\qquad g=1+(t+1)y,\qquad x=1+gu_t.
\]
Because $|y|\ge F$ and $|t+1|\ge m-1\ge7$,
\[
 |g|\ge(m-1)F-1>1.
\]
Thus $x=0$ would imply $gu_t=-1$, impossible in integers.
We have $x\ge1$ and consequently $|q|\ge2$. The literal retained
request is $q^t$. Its signed exactness forces $t\ge0$, and the
progression then gives $t\ge m$. Notice that the proof does not
replace the digit row by a statement about $B$ alone; its unrestricted
high quotient would invalidate that shortcut.

For either growth choice, positivity now yields $K\ge\ell t$.
Completeness uses the cofinal positive progression $m(2z+1)$:
choose $t$ above the finite proof length and the local free parameters,
then use the source construction with this $t$. The numerical
estimates and the order of the subsequent choices are justified below.
\end{proof}

\subsection{The estimates needed at minimum eight}

The literal square-shell constructors obey a sharper bound than
the earlier conservative $t\ge16$ presentation. For integers
$0\le r,s<T$,
\[
 E(r,s)\le8T^2-2T+1<8T^2,
\]
and the other three constructors are no larger. Thus, if both
arguments are strictly below $t^e$ and $t\ge8$, every constructor
output is strictly below $t^{2e+1}$. With
$e_0=1$ and $e_{d+1}=2e_d+1$, one obtains
$e_d=2^{d+1}-1$.

The authenticated source constructor depths are
\[
 4,\;3,\;3,\;3,\;3,\;5,\;4,\;3,\;19.
\]
Consequently every local formula label is below
$t^{2^{20}-1}$. Set $L=t^\ell$. Then $L\ge t$ and
$F=t^K\ge L^t$. No bound on the deleted sum of derived labels is
needed here.

\begin{lemma}[Boundary-valid label separation]\label{lem:label}
Let $t,L,F$ be positive integers with $L\ge t$ and $F\ge L^t$.
If two multisets of labels in $[1,L]$, with at most $t$ and
$2t+1$ entries respectively, satisfy
\[
 \prod_{j\in U}(1+jF)=\prod_{j\in V}(1+jF),
\]
then their label multiplicities agree.
\end{lemma}

\begin{proof}
If the factor counts differ, let $r$ be the smaller count.
Then $r\le t$, and comparison of the two products gives
\[
 (F+1)^{r+1}\le(1+LF)^r\le L^r(F+1)^r,
\]
contradicting $F+1>L^t$. Their common count $s$ is therefore
at most $t$. If $s=0$, both multisets are empty; assume henceforth
$s\ge1$. Each coefficient of either product polynomial
in an indeterminate lies between $1$ and
\[
 \binom{s}{k}L^k\le s^{s-k}L^k\le L^s\le F.
\]
Thus every coefficient of their difference has absolute value
strictly below $F$, even when one coefficient equals $F$.
The least nonzero coefficient of a difference vanishing at $F$
would have to be divisible by $F$, a contradiction.
The product polynomials agree. Their roots $-1/j$ recover
the multiplicities.
\end{proof}

Truth has label zero and hence weight $1+0F=1$; its factors are
omitted before applying the positive-label lemma.
The source's multiplicative Horn updates therefore recover the
same finite rule-flow balance. If falsity were outside the closure
of truth under the selected rules, every outside conclusion would
require an outside premise. Summing that balance over outside labels
would be nonpositive, whereas the target falsity contributes one.
This yields a finite C222 refutation even if the selected rules have
cycles.

Two other conservative numerical estimates have explicit replacements.
Since $K\ge\ell t$ and $t\ge8$,
\[
 F\ge8^{8\ell}>2^{64}.
\]
After the unchanged beta sign recovery, $y\ge F$ and
$x\ge2+(t+1)y$. Therefore
$2+t+x+y+F<2x$, $x>2^{32}$ and the retained local polynomial budget
gives
\[
 \mathrm{Local}<(2x)^{2^{23}}<2^x.
\]
Also
\[
 F\ge8^{2t}>16^t>
   \operatorname{lcm}(1,\ldots,t)^2.
\]
This supplies the radius for the inherited factorial-free binomial
CRT selector. No factorial-rounding lemma is being applied outside
its range.

Finally
$v_2(F)=K v_2(t)\ge\ell t$ exceeds both
$\lfloor\log_2t\rfloor$ and $v_2(t!)$. Hence
$\binom{F+t}{t}$ is odd. Multiplying $F$ by a suitable positive
odd multiple supplies a beta spacing divisible by $t!$ and larger
than all finite pending states. After choosing the beta code and
the every-prime CRT residues, choose the encoded index $C$ even
and larger than all other requested indices. Only then choose
$H,q$ and the Pell frame. The choices are cofinal and noncircular.
These facts complete the soundness and completeness assertions
in Lemma~\ref{lem:signed-horizon}.

\subsection{Two literal quartics}

The dimension-focused choice is
\[
 t=8+16z,\qquad K=1048576t.
\]
Affine elimination of $t$ gives $Q_{212}$ with 212 integer
coordinates, 151 quadratic residuals and 5,706 collected monomials.
Its exact expanded coefficient height is
\[
 316912650058804747029598371840.
\]

The balanced choice is
\[
 t=8+16z,\qquad U=2t^2,\qquad K=512U^2=2048t^4.
\]
At $t=8$ the inequality $2048t^3\ge2^{20}$ is already an equality.
This gives $Q_{214}$ with 214 integer coordinates, 153 residuals,
5,875 collected monomials and exact coefficient height 655,360.
Its maximum residual coefficient is 512.

For each profile, replace the source's horizon row by
$t-m-2mz=0$, substitute $t=m+2mz$ into every other row, and
delete this row and the coordinate $t$. This is a polynomial
bijection on zeros of the \emph{modified} source, with inverse
restoring the displayed $t$. Affine substitution keeps every
residual quadratic.

For the balanced profile, change only the two protected growth
rows to $U-2t^2$ and $K-512U^2$ before eliminating $t$.
Subtract twice the square-definition row from the frame row,
and add the square-definition row to the orientation row.
These integral row operations have explicit inverses and remove
duplicated $t^2$ terms.

The finite checker also verifies that these edits commute with
the inherited six coordinate eliminations and optimization chain.
The radical directions vanish on the protected coordinates;
the affine projection uses only the stated local constructor
rows; later charts leave the horizon and growth coordinates
unchanged; and row operations do not modify their defining rows.
Thus the emitted equations implement the stated arithmetic
source change, rather than an unsupported edit to an optimized
polynomial.

\begin{corollary}
For $r=212,214$, ZFC proves
\[
 \neg\exists\mathbf x\in\Z^r\;Q_r(\mathbf x)=0
 \quad\longleftrightarrow\quad\Con(\mathrm{ZFC}).
\]
\end{corollary}
The complete $Q_{214}$ residual list appears in
Appendix~\ref{app:quartic}; both collected polynomials are in
\texttt{quartic/}.
